\subsection{Сопряженные операторы в гильбертовом пространстве. Равенство \texorpdfstring{\(H=\overline{\operatorname{Im} A}\oplus\ker A^*\)}{H = cl Im A plus ker A*}}

\begin{definition}[Сопряженный оператор в гильбертовом пространстве]
	Пусть \(H_1,H_2\) --- гильбертовы пространства и
	\[
	A\in\mathcal{L}(H_1,H_2).
	\]
	Оператор
	\[
	A^*:H_2\to H_1
	\]
	называется \textbf{сопряженным} к \(A\), если
	\[
	(Ax,y)_{H_2}=(x,A^*y)_{H_1}
	\qquad
	\forall x\in H_1,\ \forall y\in H_2.
	\]
\end{definition}

\begin{remark}[Свойства сопряженного оператора]
	Сопряженный оператор в гильбертовом пространстве существует и единственен. Также
	выполнены свойства:
	\[
	I^*=I,
	\]
	\[
	(\alpha A+\beta B)^*
	=
	\overline{\alpha}A^*+\overline{\beta}B^*,
	\]
	\[
	(AB)^*=B^*A^*.
	\]
	Они доказываются прямой проверкой из определения.
\end{remark}

\begin{theorem}[Разложение гильбертова пространства]
	Пусть \(H\) --- гильбертово пространство над \(\mathbb{K}\), и
	\[
	A\in\mathcal{L}(H).
	\]
	Тогда
	\[
	H=\overline{\operatorname{Im}A}\oplus\ker A^*.
	\]
\end{theorem}

\begin{proof}
	Докажем разложение через ортогональное дополнение к образу.
	
	\textbf{1. Простая часть: докажем равенство \((\operatorname{Im}A)^\perp=\ker A^*\).}
	
	Пусть
	\[
	y\in(\operatorname{Im}A)^\perp.
	\]
	Это означает, что
	\[
	(Az,y)=0
	\qquad
	\forall z\in H.
	\]
	По определению сопряженного оператора
	\[
	(Az,y)=(z,A^*y).
	\]
	Значит,
	\[
	(z,A^*y)=0
	\qquad
	\forall z\in H.
	\]
	Отсюда
	\[
	A^*y=0.
	\]
	Следовательно,
	\[
	y\in\ker A^*.
	\]
	
	Обратно, если
	\[
	y\in\ker A^*,
	\]
	то
	\[
	A^*y=0.
	\]
	Тогда для любого \(z\in H\)
	\[
	(Az,y)=(z,A^*y)=(z,0)=0.
	\]
	Значит,
	\[
	y\in(\operatorname{Im}A)^\perp.
	\]
	
	Итак,
	\[
	(\operatorname{Im}A)^\perp=\ker A^*.
	\]
	
	\textbf{2. Сложная часть: применим теорему о проекции.}
	
	Знаем, что для любого множества \(S\subset H\)
	\[
	S^\perp=\left(\overline{S}\right)^\perp.
	\]
	Поэтому
	\[
	(\operatorname{Im}A)^\perp
	=
	\left(\overline{\operatorname{Im}A}\right)^\perp.
	\]
	Так как
	\[
	\overline{\operatorname{Im}A}
	\]
	--- замкнутое подпространство в \(H\), то по теореме Рисса о проекции
	\[
	H
	=
	\overline{\operatorname{Im}A}
	\oplus
	\left(\overline{\operatorname{Im}A}\right)^\perp.
	\]
	Но
	\[
	\left(\overline{\operatorname{Im}A}\right)^\perp
	=
	(\operatorname{Im}A)^\perp
	=
	\ker A^*.
	\]
	Следовательно,
	\[
	H=\overline{\operatorname{Im}A}\oplus\ker A^*.
	\]
\end{proof}

\begin{remark}
	Аналогично, применяя доказанное равенство к оператору \(A^*\), получаем
	\[
	H=\overline{\operatorname{Im}A^*}\oplus\ker A.
	\]
\end{remark}